\section{Conclusion}\label{sec:conclusion}

The right to propose a public standard and the votes needed to approve it are separate institutional instruments. In a three-country economy where firms invest in irreversible open compatibility technologies before governments decide whether to mandate a preexisting public interface, the distinction changes which private investments are supportable as pure-strategy subgame-perfect equilibria.

Our principal result is a sharp constitutional threshold on a specified region of economic primitives. If one or two affirmative votes are sufficient, granting independent proposal rights to two governments eliminates every privately retained regional-standard terminal pure equilibrium; one agenda holder cannot do so, because the excluded country may suppress a proposal that both regional founders favor. Under unanimity, regional-standard equilibria persist regardless of who can initiate a vote. The proof covers continuous heterogeneous investments by all three firms, not simply an assumed symmetric regional technology.

The same institutional reform has no unconditional efficiency ranking. The nonmember's national veto, the world-welfare effect of retrofitting after private expenditure is sunk, and the welfare ranking of alternative initial public and private investment paths have three distinct cutoffs. Their ordered difference shows why desirable public standardization at the outset can coexist with an undesirable retrofit after private investment. A rule that produces a unique public equilibrium can also remove a privately developed equilibrium with higher ex-ante world welfare.

These conclusions rest on a tightly specified public fallback, a no-transfer voting game, and a tractable compatibility technology. They do not settle standardization with licensing, patent ownership, endogenous public design, or mixed strategies. What the model isolates is a mechanism in which \emph{constitutional proposal rights, rather than the voting quota alone, govern whether firms find it worthwhile to build a regional technical alternative}. The welfare implications of that mechanism depend on when regulation intervenes relative to sunk technological commitments.
